\documentclass[11pt, a4paper]{article}
\usepackage[margin=2.5cm]{geometry}
\usepackage{booktabs}
\usepackage{array}
\usepackage{longtable}
\usepackage{xcolor}
\usepackage{hyperref}
\usepackage{microtype}
\usepackage{parskip}
\usepackage{enumitem}
\usepackage[T1]{fontenc}
\usepackage[utf8]{inputenc}

% Colour definitions for priority levels
\definecolor{highprio}{HTML}{CC0000}
\definecolor{medprio}{HTML}{CC8800}
\definecolor{lowprio}{HTML}{007700}

\hypersetup{
  colorlinks = true,
  linkcolor  = blue!60!black,
  urlcolor   = blue!60!black,
}

\title{\textbf{ERBOT Results Summary}\\[4pt]
       \large NCVR 5-Party Voter Registration Dataset}
\author{ERBOT v0.2.0}
\date{Generated: 2026-03-25 \quad|\quad Small run: $<$1 hr \quad|\quad Large run: $\approx$45 hrs}

\begin{document}
\maketitle
\tableofcontents
\vspace{1em}
\hrule
\vspace{1em}

% ─────────────────────────────────────────────────────────────────────────────
\section{Dataset Overview}

\begin{table}[h!]
\centering
\begin{tabular}{ll}
\toprule
\textbf{Property} & \textbf{Value} \\
\midrule
Dataset      & NCVR 5-Party Voter Registration (ocp=20, modrec=2) \\
Source files & \texttt{ncvr\_numrec\_1000000\_modrec\_2\_ocp\_20\_myp\_\{0,1\}\_nump\_5.csv} \\
Parties used & 2 (party 0 and party 1 from the 5-party collection) \\
Fields       & 4 content fields: \texttt{givenname}, \texttt{surname}, \texttt{suburb}, \texttt{postcode} \\
Task         & Record linkage (two-source) \\
Overlap rate & 20\% (\texttt{ocp=20}); modification rate: 2 fields per duplicate (\texttt{modrec=2}) \\
True matches & $\approx$880 cross-party pairs (varies slightly by sample) \\
\bottomrule
\end{tabular}
\end{table}

The NCVR 5-Party benchmark is derived from the North Carolina Voter Registration
database. The generator produces $N$ records per party with a controlled overlap
percentage and modification rate. Records that appear in both parties share the
same \texttt{recid} but may differ in up to \texttt{modrec} fields. This benchmark
tests cross-party record linkage under realistic personal-name and address noise.
Two runs were conducted: a \textbf{small run} (5K records per party) and a
\textbf{large run} (50K records per party) to assess scalability.

% ─────────────────────────────────────────────────────────────────────────────
\section{Pipeline Configuration}

\begin{table}[h!]
\centering
\begin{tabular}{llll}
\toprule
\textbf{Stage} & \textbf{Setting} & \textbf{Small (5K/party)} & \textbf{Large (50K/party)} \\
\midrule
Input       & Records total       & 10,000      & 100,000 \\
Input       & True matches        & $\approx$880 & $\approx$880 \\
Input       & Match density       & $\sim$8.8\% & $\sim$0.88\% \\
Blocking    & Method              & Prefix-3 on \texttt{surname} & Prefix-3 on \texttt{surname} \\
Blocking    & Candidate pairs     & $\sim$10,600 (est.) & 1,059,972 (truncated from 2,121,188) \\
Blocking    & Reduction ratio     & $\approx$1.000 & 1.000 \\
Similarity  & Fields computed     & 4 & 4 \\
Weights     & Method              & Equal (0.25 each) & Equal (0.25 each) \\
Clustering  & Methods run         & All 9 (incl.\ SVM) & 8 (SVM excluded; hclust/pam/gc failed) \\
Merge       & Strategy            & Best & Best \\
Final       & Clusters            & 1,910 & 94,706 \\
\bottomrule
\end{tabular}
\end{table}

\textbf{Blocking note:} Prefix-3 on \texttt{surname} achieves near-total pair
reduction (RR $\approx$ 1.000 due to rounding at this scale). In the large run
the blocker generated 2,121,188 raw pairs which exceeded the 2M pair budget and
was truncated to approximately 2M; the final sparse matrix has 2,119,944
non-zero entries. This truncation may silently drop true match pairs and
directly suppresses recall.

\textbf{Match density note:} Both runs target the same absolute number of true
matches ($\approx$880) because the ground truth is defined by recid overlap in
the sampled subsets. The large run has 10$\times$ more records but the same
matches, giving a 10$\times$ lower match density (0.88\% vs.\ 8.8\%). This
is the primary driver of the performance gap between small and large runs.

% ─────────────────────────────────────────────────────────────────────────────
\section{Performance Results --- Small Run (5\,000 records/party)}

\begin{table}[h!]
\centering
\small
\setlength{\tabcolsep}{4pt}
\begin{tabular}{lrrrrrrrrrrrr}
\toprule
\textbf{Method} & \textbf{ARI} & \textbf{NMI} & \textbf{VI} &
\textbf{B$^3$-P} & \textbf{B$^3$-R} & \textbf{B$^3$-F} &
\textbf{Hom.} & \textbf{Com.} & \textbf{V} &
\textbf{PF-P} & \textbf{PF-R} & \textbf{PF-F} \\
\midrule
\textbf{threshold\_cc} & \textbf{0.917} & \textbf{0.975} & \textbf{0.143} & 1.000 & 0.929 & \textbf{0.963} & 1.000 & 0.952 & \textbf{0.975} & 1.000 & 0.857 & \textbf{0.923} \\
\textbf{louvain}       & \textbf{0.917} & \textbf{0.975} & \textbf{0.143} & 1.000 & 0.929 & \textbf{0.963} & 1.000 & 0.952 & \textbf{0.975} & 1.000 & 0.857 & \textbf{0.923} \\
\textbf{label\_prop}   & \textbf{0.917} & \textbf{0.975} & \textbf{0.143} & 1.000 & 0.929 & \textbf{0.963} & 1.000 & 0.952 & \textbf{0.975} & 1.000 & 0.857 & \textbf{0.923} \\
\textbf{svm}           & \textbf{0.917} & \textbf{0.975} & \textbf{0.143} & 1.000 & 0.929 & \textbf{0.963} & 1.000 & 0.952 & \textbf{0.975} & 1.000 & 0.857 & \textbf{0.923} \\
hclust\_avg            & 0.276 & 0.713 & 1.315 & 0.537 & 0.929 & 0.681 & 0.583 & 0.920 & 0.713 & 0.231 & 0.857 & 0.364 \\
pam                    & 0.158 & 0.672 & 1.601 & 0.537 & 0.786 & 0.638 & 0.583 & 0.792 & 0.672 & 0.167 & 0.571 & 0.258 \\
leiden                 & 0.000 & 0.849 & 1.000 & 1.000 & 0.500 & 0.667 & 1.000 & 0.737 & 0.849 & 0.000 & 0.000 & 0.000 \\
hclust\_ward           & 0.068 & 0.542 & 1.944 & 0.429 & 0.857 & 0.571 & 0.409 & 0.801 & 0.542 & 0.111 & 0.714 & 0.192 \\
gc                     & 0.000 & 0.000 & 2.807 & 0.143 & 1.000 & 0.250 & 0.000 & 1.000 & 0.000 & 0.077 & 1.000 & 0.143 \\
\midrule
\textit{final}         & 0.917 & 0.975 & 0.143 & 1.000 & 0.929 & 0.963 & 1.000 & 0.952 & 0.975 & 1.000 & 0.857 & 0.923 \\
\bottomrule
\end{tabular}
\end{table}

% ─────────────────────────────────────────────────────────────────────────────
\section{Performance Results --- Large Run (50\,000 records/party)}

\begin{table}[h!]
\centering
\small
\setlength{\tabcolsep}{4pt}
\begin{tabular}{lrrrrrrrrrrrr}
\toprule
\textbf{Method} & \textbf{ARI} & \textbf{NMI} & \textbf{VI} &
\textbf{B$^3$-P} & \textbf{B$^3$-R} & \textbf{B$^3$-F} &
\textbf{Hom.} & \textbf{Com.} & \textbf{V} &
\textbf{PF-P} & \textbf{PF-R} & \textbf{PF-F} \\
\midrule
\textbf{threshold\_cc} & \textbf{0.195} & \textbf{0.956} & \textbf{0.908} & \textbf{0.991} & 0.560 & \textbf{0.715} & \textbf{0.997} & 0.917 & \textbf{0.956} & \textbf{0.533} & 0.119 & \textbf{0.195} \\
louvain                & 0.133 & 0.947 & 1.073 & 0.933 & 0.560 & 0.700 & 0.980 & 0.916 & 0.947 & 0.152 & 0.119 & 0.134 \\
label\_prop            & 0.133 & 0.947 & 1.073 & 0.933 & 0.560 & 0.700 & 0.980 & 0.916 & 0.947 & 0.152 & 0.119 & 0.134 \\
leiden                 & 0.000 & 0.951 & 1.000 & 1.000 & 0.500 & 0.667 & 1.000 & 0.907 & 0.951 & 0.000 & 0.000 & 0.000 \\
hclust\_avg$^*$        & 0.000 & 0.000 & 9.781 & 0.001 & 1.000 & 0.002 & 0.000 & 1.000 & 0.000 & 0.001 & 1.000 & 0.001 \\
hclust\_ward$^*$       & 0.000 & 0.000 & 9.781 & 0.001 & 1.000 & 0.002 & 0.000 & 1.000 & 0.000 & 0.001 & 1.000 & 0.001 \\
pam$^*$                & 0.000 & 0.000 & 9.781 & 0.001 & 1.000 & 0.002 & 0.000 & 1.000 & 0.000 & 0.001 & 1.000 & 0.001 \\
gc$^*$                 & 0.000 & 0.000 & 9.781 & 0.001 & 1.000 & 0.002 & 0.000 & 1.000 & 0.000 & 0.001 & 1.000 & 0.001 \\
\midrule
\textit{final}         & 0.195 & 0.956 & 0.908 & 0.991 & 0.560 & 0.715 & 0.997 & 0.917 & 0.956 & 0.533 & 0.119 & 0.195 \\
\bottomrule
\end{tabular}
\end{table}

$^*$ \textit{Method failed at runtime and fell back to degenerate singleton assignment.
hclust\_avg/ward and pam exceeded R's 65,536-observation distance-matrix limit;
gc attempted a 74.5\,GiB dense allocation. These methods are now excluded from
\texttt{run\_ncvr5\_large.R}.}

\textit{Metrics: ARI = Adjusted Rand Index; NMI = Normalized Mutual Information;
VI = Variation of Information (lower is better); B$^3$ = B-cubed; V = V-measure;
PF = Pairwise F-score; Hom.\ = Homogeneity; Com.\ = Completeness.}

% ─────────────────────────────────────────────────────────────────────────────
\section{Scale Comparison}

\begin{table}[h!]
\centering
\begin{tabular}{lrrl}
\toprule
\textbf{Metric} & \textbf{Small (5K/party)} & \textbf{Large (50K/party)} & \textbf{Comment} \\
\midrule
Records            & 10,000    & 100,000    & 10$\times$ scale-up \\
True matches       & $\sim$880 & $\sim$880  & Same absolute count \\
Match density      & 8.8\%     & 0.88\%     & 10$\times$ sparser \\
Best ARI           & 0.917     & 0.195      & $-$0.72 drop \\
Best B$^3$-F       & 0.963     & 0.715      & $-$0.25 drop \\
Best method        & threshold\_cc / louvain / label\_prop / svm (tied) & threshold\_cc & Consistent winner \\
Final correct      & Yes       & Yes        & Merge logic working \\
Final clusters     & 1,910     & 94,706     & Most records in singletons at large scale \\
Blocking RR        & $\approx$1.000 & 1.000 & Identical (rounding artefact) \\
B$^3$ Precision    & 1.000     & 0.991      & Maintained well \\
B$^3$ Recall       & 0.929     & 0.560      & Major recall drop \\
Runtime            & $<$1 hr   & $\approx$45 hrs & hclust/pam/gc failures inflated time \\
\bottomrule
\end{tabular}
\end{table}

The performance collapse from small to large is driven almost entirely by
\textbf{recall}: precision stays high (1.0 $\to$ 0.99) but recall drops from
0.93 to 0.56. At 50K records per party, the sparse surname prefix blocks contain
so many non-match pairs that true match pairs are diluted, and some are
truncated by the pair budget cap. This is a blocking problem, not a clustering
problem.

% ─────────────────────────────────────────────────────────────────────────────
\section{Final Cluster Assignment}

In both runs the \texttt{final} merge strategy correctly selected
\textbf{threshold\_cc} (highest ARI).

\begin{table}[h!]
\centering
\begin{tabular}{lrr}
\toprule
\textbf{Statistic} & \textbf{Small} & \textbf{Large} \\
\midrule
Total clusters   & 1,910  & 94,706 \\
Records          & 10,000 & 100,000 \\
Mean cluster size & 5.2   & 1.06 \\
\bottomrule
\end{tabular}
\end{table}

The large run's 94,706 clusters from 100,000 records (mean size $\approx$1.06)
indicates that nearly all records are assigned as singletons --- correctly
reflecting the low match density (only 880 true matches exist in 100K records).
The small run's 1,910 clusters from 10,000 records (mean 5.2) is more
compact but reflects the much higher match density at that scale.

% ─────────────────────────────────────────────────────────────────────────────
\section{Method-by-Method Interpretation}

\subsection*{Threshold-CC (Best overall --- ARI 0.917 / 0.195)}
The consistent winner at both scales. At small scale it ties with louvain,
label\_prop, and svm at ARI = 0.917 and perfect precision (B$^3$-P = 1.0).
At large scale it retains near-perfect precision (0.991) but recall falls to
0.560. The connected-components approach is well-suited to the NCVR similarity
graph: at a threshold of 0.5 the graph separates cleanly, and the four-field
equal-weight profile creates a well-calibrated similarity score for voter
records with controlled 2-field modifications.

\subsection*{Louvain / Label Propagation (Tied-best small; second large)}
Identical to threshold-CC at small scale; slightly weaker at large scale
(ARI 0.133 vs.\ 0.195). This divergence suggests that at low match density
the community structure in the graph is less clean, and Louvain/LP begin to
over-merge some blocks. The lower PairF-P (0.152 vs.\ 0.533 for threshold-CC)
at large scale confirms more false-positive links.

\subsection*{SVM (Best small; not run large)}
Ties at ARI = 0.917 in the small run. SVM is excluded from the large run
script (correctly --- it is too slow at 100K scale). At 5K/party scale it
adds no benefit over the threshold-CC family on this structured dataset, but
it would be worth including in any small-scale parameter sweep.

\subsection*{Leiden (Degenerate at both scales --- ARI = 0.000)}
Perfect precision (1.0) and completeness (0.74/0.91) but ARI = 0 and PairF = 0
at both scales. Leiden over-partitions: it splits true duplicate pairs into
separate singleton clusters rather than merging them. The resolution parameter
is too high for NCVR's similarity graph topology. Needs tuning.

\subsection*{hclust\_avg (Partial failure)}
Works at small scale (ARI = 0.276, second-best after the graph quartet) but
fails at large scale with R's 65,536 observation limit. Even at small scale
it underperforms graph methods substantially.

\subsection*{hclust\_ward / PAM / GC (Fail or degenerate)}
hclust\_ward and PAM fail at large scale (65K limit). GC degenerates at both
scales (ARI = 0): at small scale it produces a near-singleton partition with
perfect recall (B$^3$-R = 1.0) but near-zero precision (0.143); at large scale
it crashes with a 74.5\,GiB allocation attempt. GC is not viable for NCVR.

% ─────────────────────────────────────────────────────────────────────────────
\section{Key Findings}

\begin{enumerate}[leftmargin=*, label=\arabic*.]
  \item \textbf{Threshold-CC is the best and most scalable method.} It wins
        at both scales, is fast, and avoids the memory failures that affect
        distance-matrix methods. It is the correct default for NCVR-style
        4-field voter linkage.

  \item \textbf{The performance drop from small to large is a blocking/density
        problem, not a clustering problem.} Precision is maintained (0.991) but
        recall drops from 0.929 to 0.560. At 50K/party, true match pairs are
        sparser and some are truncated by the 2M pair budget cap. The clustering
        algorithm cannot recover pairs that blocking discards.

  \item \textbf{Pair budget truncation is active at large scale.}
        The blocker found 2,121,188 raw pairs ($>$2M budget) and truncated to
        $\approx$2M. Some true match pairs may have been silently dropped at this
        step. The budget should be raised or a smarter blocking strategy adopted.

  \item \textbf{hclust\_avg/ward, PAM, and GC are not viable at 100K scale.}
        All three fail with R's 65K distance-matrix limit or memory exhaustion.
        They have been removed from the large run script. They should also be
        flagged as ``small-scale only'' in documentation.

  \item \textbf{Equal weights work well for NCVR.} With 4 equally noisy fields
        (name, surname, suburb, postcode), equal weighting is a reasonable
        choice. Unlike CORA (15 fields of varying quality), there is no evidence
        that weight learning would substantially improve results here.

  \item \textbf{The \texttt{final} merge strategy is working correctly.}
        It selected threshold-CC (highest ARI) in both runs. The earlier merge
        bug seen in CORA has been resolved.

  \item \textbf{Leiden consistently degenerates.} ARI = 0 at both scales
        across NCVR and CORA suggests the default resolution parameter is
        unsuitable for these datasets. Leiden should not be included in
        production runs without prior resolution tuning.
\end{enumerate}

% ─────────────────────────────────────────────────────────────────────────────
\section{Recommendations for Next Steps}

\begin{table}[h!]
\centering
\begin{tabular}{cp{0.78\linewidth}}
\toprule
\textbf{Priority} & \textbf{Action} \\
\midrule
{\color{highprio}\textbf{High}}   & Increase pair budget beyond 2M (or make it adaptive) to avoid silent truncation of true matches at large scale. \\[4pt]
{\color{highprio}\textbf{High}}   & Add pair completeness measurement at Stage 3 to quantify how many true matches the blocker discards. Without this, recall numbers are lower bounds, not true estimates. \\[4pt]
{\color{highprio}\textbf{High}}   & Remove leiden and gc from default method lists (or add a degenerate-detection guard) --- both produce ARI = 0 consistently without tuning. \\[4pt]
{\color{medprio}\textbf{Medium}}  & Tune Leiden resolution parameter (try 0.05--0.2) on the small run before running large-scale. \\[4pt]
{\color{medprio}\textbf{Medium}}  & Try more selective blocking (e.g.\ phonetic key on \texttt{surname} + exact \texttt{postcode}) to improve recall without exploding pair count. \\[4pt]
{\color{medprio}\textbf{Medium}}  & Run all 5 parties (not just p0 and p1) to assess multi-party linkage performance. \\[4pt]
{\color{lowprio}\textbf{Low}}     & Profile whether the 45-hour large run time is dominated by the failed hclust/pam/gc attempts (expected: yes). Re-time with fixed method list. \\[4pt]
{\color{lowprio}\textbf{Low}}     & Experiment with field-specific similarity (Jaro-Winkler for names, exact match for postcode) rather than auto-detected string similarity. \\
\bottomrule
\end{tabular}
\end{table}

% ─────────────────────────────────────────────────────────────────────────────
\section{Comparison Across Datasets}

\begin{table}[h!]
\centering
\begin{tabular}{lllp{0.35\linewidth}}
\toprule
\textbf{Metric} & \textbf{CORA} & \textbf{Affiliation} & \textbf{NCVR5 (small / large)} \\
\midrule
Records        & 1,879  & 2,260  & 10,000 / 100,000 \\
Fields         & 15     & 1      & 4 \\
Task           & Dedup  & Dedup  & Linkage \\
Best ARI       & 0.543  & 0.076  & 0.917 / 0.195 \\
Best B$^3$-F   & 0.714  & 0.460  & 0.963 / 0.715 \\
Best method    & SVM    & threshold\_cc & threshold\_cc \\
Final correct  & No (bug) & Yes  & Yes / Yes \\
Main bottleneck & Merge bug, equal weights & Blocking, 1 field & Blocking recall, pair budget \\
\bottomrule
\end{tabular}
\end{table}

NCVR5 at small scale is the strongest result across all three datasets
(ARI = 0.917), reflecting a well-structured 4-field linkage task with controlled
noise. At large scale the density effect degrades performance substantially.
The affiliation dataset remains the hardest challenge (ARI = 0.076), driven
by the single-field limitation.

% ─────────────────────────────────────────────────────────────────────────────
\section{Output Files}

\begin{table}[h!]
\centering
\begin{tabular}{ll}
\toprule
\textbf{File} & \textbf{Description} \\
\midrule
\texttt{small/er\_performance.csv}        & Full metric table --- small run (10 methods $\times$ 13 metrics) \\
\texttt{small/er\_predictions.csv}        & Final cluster assignments --- small run (10,000 records) \\
\texttt{small/plot\_method\_comparison.pdf} & Grouped bar chart: ARI, B$^3$-F, NMI per method (small) \\
\texttt{small/plot\_performance\_heatmap.pdf} & Colour heatmap: all metrics $\times$ methods (small) \\
\texttt{small/plot\_cluster\_sizes.pdf}   & Cluster size distribution histogram (small) \\
\texttt{small/plot\_precision\_recall.pdf} & Precision vs.\ Recall scatter (small) \\
\texttt{large/er\_performance.csv}        & Full metric table --- large run (9 methods $\times$ 13 metrics) \\
\texttt{large/er\_predictions.csv}        & Final cluster assignments --- large run (100,000 records) \\
\texttt{large/run\_large.log}             & Full console log of large run \\
\texttt{large/plot\_method\_comparison.pdf} & Grouped bar chart (large) \\
\texttt{large/plot\_performance\_heatmap.pdf} & Colour heatmap (large) \\
\texttt{large/plot\_cluster\_sizes.pdf}   & Cluster size distribution histogram (large) \\
\texttt{large/plot\_precision\_recall.pdf} & Precision vs.\ Recall scatter (large) \\
\texttt{ncvr5\_data\_profile.pdf}         & Dataset profile: field types, distributions, missingness \\
\texttt{NCVR5\_RESULTS\_SUMMARY.pdf}      & This document \\
\bottomrule
\end{tabular}
\end{table}

\end{document}
